\section{Introduction}\label{sec:introduction}
A current economic choice changes a future that is costly to evaluate. A
household chooses consumption while investing in its preferences; a firm
adjusts capital against the policies of its rivals; an economist recomputes
these choices under alternative charges or constraints. Neural Bellman
Operators (NBO) evaluate a specified continuation policy and use a feasible
actor to improve current decisions. The question is when paying for a learned
continuation makes these decisions more accurate, or less expensive at a
specified economic accuracy, than directly simulating or approximating them.

This question involves three distinct objects. The continuation equation
specifies the future being valued. The actor specifies the feasible decision
actually taken. The economic error account specifies what is known about its
payoff. A small regression loss does not, by itself, answer the economic
question; neither does a short training clock. The comparison must charge
construction, all action queries, unsuccessful accuracy checks, and the
verification required to return a decision. It must also allow competing
procedures to reuse the information that their decision problems permit.

The paper develops this evaluation--decision--payoff chain for controlled
economies. Its original differential and monotone formulations, recursive
utility, endogenous preferences, temporal selves, and strategic interaction
remain substantive components. Policy evaluation is always policy specific.
A utility-domain restriction cannot be discarded when the approximation is
reused, and a best response against fixed rivals is not a proof of equilibrium.
These distinctions determine which continuation can be shared and when it
must be refreshed.

Our first contribution is an explicit connection between continuation reuse
and decision loss. A centered error, rather than an arbitrary level error,
controls the loss from an approximate maximizing action. A transport account
then bounds the additional loss when future economic primitives change.
The resulting reuse region includes bounded changes in transitions, rewards,
and recursive evaluation operators; it is not confined to tasks with an
identical future. The bound makes the relevant stability and utility-domain
conditions explicit. It also explains why prediction-risk rankings need not
coincide with the payoff rankings of competing policies.

Our second contribution is a work-to-accuracy implementation of this account.
For a finite-support instance of the nonlinear capital economy, an outward
interval calculation certifies the true objective's derivative and a uniform
strong-concavity bound. Every returned action is assessed against the full
continuous scalar action interval, not merely a computed candidate list.
A frozen sequence of fitting budgets stops at the first successful economic
check and charges every attempt. This supplies a realized stopping cost,
whereas the earlier simultaneous catalogue closure supplies a valid but
retrospective accuracy comparison. Both calculations remain in the paper.

The completed design varies the number of decisions from one to 1,024 in
dimensions ten and fifty. Temporary withdrawal charges, current utility
weights, adjustment costs, and caps generate the decisions while the future
reference policy stays fixed. Seven complete procedures include NBO,
conventional quadratic and radial-basis continuation regressions, a shared
task-conditioned actor, cached sample-average optimization, full-support
enumeration, and a current-reward rule. Every dimension--volume cell is
executed under three fixed fitting and cache streams. The integrated execution
replays this already frozen design; it does not create additional independent
streams by repeating a computation. The resulting economic certificates are
deterministic for the stored finite economy.

NBO reaches the prescribed scalar-interval tolerance in all 36 of its
complete runs. This identifies successful learned-continuation decisions, not
uniform neural superiority. The non-neural surrogates and simulation-based
procedures also achieve the target and remain on the measured cost frontier.
The shared actor and current-only procedure retain their unsuccessful checks.
Common exogenous actions and each predictor's own selected actions are both
assessed against full-support continuation means; caches are regenerated
across the three streams. Thus absolute risk, centered risk, and actual
finite-menu decision loss can be compared without treating simulation noise
as the target value or selecting the diagnostic action pair in favor of one
predictor.

This evidence complements the original capital experiments rather than
replacing them. In the previously unpiloted fifty-dimensional Long and
Intermediate menus, the protected NBO payoff advantages over materialized
Raw and direct-policy actors exceed $10^{-4}$. The earlier conditional
continuation-risk assessment yields eight comparisons favoring NBO, four
favoring Raw, and twelve unresolved comparisons. None of its eight
64-path comparisons favors NBO. Exact closure of the original simultaneous
payoff intervals makes final NBO the least-cost \emph{certified} candidate in
one of eight cells, while cheaper unresolved candidates remain. Those
statements retain their original finite populations, probability allocations,
and chronology; they are not restated as results of the new stopping design.

The continuous-time evidence also keeps its own economic target. Independent
confirmation favors NBO over the original neural-HJB implementation in both
dimensions. Against the strengthened exact-trace HJB procedure, only the
low-volatility fifty-dimensional cell establishes the declared margin against
both deployments. The signed Bellman assessment has positive continuous-economy
gain bounds, while its signed continuation-correction intervals contain zero.
Earlier unresolved NBO--Raw comparisons, the cheaper Raw clock, and
nonpositive mechanism bounds are preserved. A finite scalar certificate does
not remove diffusion, implementation, information, or full-control errors
from these separate accounts.

The principal economic exercise in the new design is comparative statics
for temporary resource charges. Its primary target is foregone transformed
welfare at the chosen action. We derive the sign of the optimal charge
response and translate the target into an external log-utility compensation.
This is a specified counterfactual exercise in the paper's capital economy,
not an estimate of an observed tax reform. The task distribution defines the
exercise and is not presented as an estimated population of policy requests.
The results show when a common continuation supports these decisions and
what its construction and certification cost. They do not require that a
neural representation be the best approximation for every economic model.


The final link is from current decisions to policy accuracy. A composition
theorem propagates certificates for the implemented policy's own continuation
through monotone evaluation operators, with explicit state-domain coverage.
A companion allocation result distributes fitting allowances across dates
at a fixed policy target. The continuous-economy class and value-transfer
errors remain explicit, and recursive utility and strategic deviations retain
their model-specific restrictions.

A separately frozen extension changes future withdrawal, technology, and
payoff weights rather than only current tasks. In its 168 services and 840
future-specific outcomes, checking and refreshing NBO certifies every
future; unchanged continuation reuse leaves 48 outcomes uncertified. The
returned neural actions also certify negative current-withdrawal responses
to stronger future production and greater future valuation. These results
identify a useful learned continuation beyond an identical-future menu.
They do not suppress the finding that the registered adaptive neural rule
costs more than fresh low-budget neural fitting in every volume cell, or the
competitive non-neural procedures. The current and future-change experiments
have distinct source freezes and streams, both preserved in the record.

The main article presents the common decision problem, its economic error
chain, and the complete primary comparisons. The technical supplement gives
full proofs, the detailed inherited capital analysis, and all computational
records. The Economic Applications companion retains the complete utility,
preference, temporal-self, and game models. This organization separates the
central argument from its derivations without changing the paper's subject
or discarding unfavorable evidence.
